\documentclass{article}
	\usepackage[UTF8]{ctex}
	\usepackage{amsmath}
	\usepackage{graphicx}
	\usepackage{underscore}
	\usepackage{geometry}
	\usepackage{anysize}
	\usepackage{amssymb}
	\date{}
	\title{最长子序列 \\ [2ex] \begin{large} 
	时间限制：C/C++ 1秒，其他语言2秒 \\
	空间限制：C/C++ 256.000MB，其他语言512.000MB 	\\
	64bit IO Format: \%lld
 			\end{large} }
	\geometry{top=10mm,bottom=15mm,left=25mm,right=25mm}
	
	\begin{document}
		\maketitle \vspace{-50pt}
		
		\section*{题目描述}
			\indent 在一个给定的数值序列$a_{1},a_{2},\dots,a_{n}$中，寻找子序列$a_{b_{1}},a_{b_{2}},\dots,a_{b_{m}}$($1 \le b_{1}<b_{2}<\dots<b_{m} \le n$)，对于$\forall i \in [1,m) , i \in \mathbb{N}$，满足以下条件：\\
\begin{equation}\nonumber
	\left\{
	\begin{aligned}
		b_{i+1}-b_{i} \ge \lceil \frac{x}{b_{i+1}} \rceil \\
		a_{b_{i+1}}-a_{b_{i}} > \lfloor \frac{b_{i+1}-b_{i}}{y} \rfloor
	\end{aligned}
	\right.
\end{equation}

	\indent 其中$\forall i \in [1,m) , i \in \mathbb{N}$的意思为对于大于等于1小于m的整数i。$\lceil s \rceil$为向上取整，结果为大于等于s的最小整数。$\lfloor  s \rfloor$为向下取整，结果为小于等于s的最大整数。\\
			\indent 求满足条件的子序列的最大长度。
		
		\section*{输入描述}
			\indent 第一行输入T($1 \le T \le 10$)，代表数据组数。\\
			\indent 对于每组数据，第一行输入n($1 \le n \le 10000$)，x($n \le x \le 5*n$)，y($1 \le y \le 5$)，第二行输入n个整数($-1000000 \le a_{j} \le 1000000$，$ 1 \le j \le n$)，代表数值序列。
			
		\section*{输出描述}
			对于每组数据，输出一个整数，为满足条件的子序列的最大长度。
		
		\section*{示例1}
			\subsection*{输入}
\noindent
2\\
1 1 1\\
1\\
5 5 1\\
1 4 3 9 16
			\subsection*{输出}
\noindent
1\\
3

		\section*{备注}	
\noindent 对于第1组数据，1为满足条件的最长子序列。
对于第2组数据，1 9 16和4 9 16皆为满足条件的最长子序列。
	\end
	{document}
	